\chapter{Desenvolvimento do Trabalho} \label{cap:desenvolvimento}

\section{Tecnologias Utilizadas}
\label{sec:tecnologias}

A etapa estática é construída sobre a infraestrutura do compilador LLVM
\cite{llvm}. Os programas são compilados com o Clang para a representação
intermediária, com \texttt{-g}, para que os metadados de depuração preservem o
vínculo entre cada instrução e a linha do código-fonte, e com \texttt{-O0},
para que nenhum acesso à memória seja eliminado por otimização. Programas com
vários arquivos são compilados separadamente e ligados num só módulo com o
\texttt{llvm-link}. Ligar a representação intermediária, em vez de concatenar
código-fonte, preserva a origem de cada instrução, resolve símbolos
\texttt{static} de mesmo nome em arquivos diferentes e respeita as definições
de pré-processador de cada um.

A análise de ponteiros usa a biblioteca SVF \cite{svf}, que trabalha sobre essa
mesma representação e oferece as consultas do tipo \textit{may} de que a
Seção~\ref{sec:tecnicas} trata. A verificação de viabilidade de caminho usa o
solucionador Z3 \cite{z3}, seguindo o uso que \citeonline{intrace} fazem dele.
A reverificação recorre a um verificador de modelos limitado, o CBMC
\cite{cbmc} ou o próprio BMC4AV, que é a única das quatro ferramentas revisadas
cujo código-fonte está disponível.

A etapa de triagem é escrita em Python, com um cliente que isola o provedor do
modelo atrás de uma interface própria, como exige o RNF5. As respostas do
modelo são pedidas em formato estruturado e validadas contra um esquema antes
de serem usadas, de modo que uma resposta malformada falha de imediato em vez
de contaminar uma medição. O cache de respostas é indexado pelo candidato, pela
configuração de \textit{prompt} e pelo modelo, o que permite recalcular uma
tabela sem repetir chamadas.

\section{Projeto e Implementação}
\label{sec:projeto}

\subsection{Visão geral}

A ferramenta percorre o ciclo detectar, contextualizar, triar, reparar e
reverificar, ilustrado na Figura~\ref{fig:pipeline}. A etapa estática produz os
candidatos e o contexto de cada um; a etapa com modelo de linguagem os ordena,
explica e propõe correções; e a própria análise estática, com apoio de um
verificador limitado, confirma as correções.

\begin{figure}[htb]
  \centering
  \caption{Pipeline da ferramenta proposta}
  \label{fig:pipeline}
  \begin{tikzpicture}[
      node distance=5mm and 4.5mm,
      etapa/.style={draw=black, rounded corners=2pt, align=center,
        minimum height=11mm, text width=29.5mm, font=\small, fill=white},
      llm/.style={etapa, dashed},
      seta/.style={-{Stealth[length=2.2mm]}, thick, color=black},
      nota/.style={font=\scriptsize, color=black, align=center}]
    \node[etapa] (cfg) {Código C +\\configuração};
    \node[etapa, right=of cfg] (e1) {Estágio 1\\candidatos};
    \node[etapa, right=of e1] (e2) {Estágio 2\\prioridade e\\mascaramento};
    \node[etapa, right=of e2] (e3) {Estágio 3\\viabilidade (Z3)};
    \node[llm, below=12mm of e3] (tri) {Triagem\\(ordena e explica)};
    \node[llm, left=of tri] (rep) {Reparo\\com testemunha};
    \node[etapa, left=of rep] (rev) {Reverificação};
    \node[etapa, left=of rev] (rel) {Relatório\\ordenado};
    \draw[seta] (cfg) -- (e1);
    \draw[seta] (e1) -- (e2);
    \draw[seta] (e2) -- (e3);
    \draw[seta] (e3) -- node[right, nota] {registro de\\contexto} (tri);
    \draw[seta] (tri) -- (rep);
    \draw[seta] (rep) -- (rev);
    \draw[seta] (rev) -- (rel);
    \draw[seta, dashed] (rev.north) to[bend left=18] node[left, nota, pos=0.35] {nova\\análise} (e1.south);
    \node[nota, above=1mm of e2] {só descartam com prova de impossibilidade};
    \node[nota, below=1mm of rep] {etapas com \llm\ (tracejadas): nunca descartam};
  \end{tikzpicture}
  \fonte{Elaborada pelos autores.}
\end{figure}

\subsection{Entrada e configuração}

O modelo de interrupções vem de um arquivo de configuração escrito à mão, que
nomeia os pontos de entrada de cada fluxo, suas prioridades e as primitivas de
mascaramento da plataforma. Inferir esse modelo com um \llm, como
\citeonline{iris} fazem com especificações de bibliotecas, foi considerado e
descartado. As primitivas de mascaramento formam um conjunto fechado por
plataforma, os dois \textit{benchmarks} já distribuem a configuração, e
encontrar os pontos de entrada, que é a parte crítica para o recall, é uma
busca textual pela API de registro de interrupções, não um julgamento. O
arquivo tem ainda um efeito colateral útil sobre o custo, já que a análise
percorre apenas o código alcançável a partir dos fluxos declarados, e não o
repositório inteiro.

\subsection{Estágio 1, derivação de candidatos}

O primeiro estágio segue o desenho de \citeonline{intrace} e ignora de
propósito condições de desvio e estado das interrupções, para não perder nada.
Ele identifica os fluxos a partir da configuração, tratando cada ponto de
entrada como raiz de uma análise de alcançabilidade; identifica as posições
compartilhadas por análise \textit{may}, sem se restringir a variáveis globais,
já que uma variável local cujo endereço circula em ponteiros globais também é
compartilhada; enumera os acessos a essas posições com todo o contexto exigido
pelo RF3; e deriva os candidatos.

Os candidatos são as combinações $(A_1, B, A_2)$ que se encaixam em algum dos
quatro padrões não serializáveis do Quadro~\ref{quadro:padroes}, com $A_1$ e
$A_2$ no fluxo interrompido, nessa ordem, e $B$ num fluxo capaz de interrompê-lo
entre os dois. A ordem entre $A_1$ e $A_2$ vem do grafo de fluxo de controle da
função que os contém.

Um detalhe sobre laços é decisivo para o recall. Um único acesso no
código-fonte, dentro de um laço que pode repetir, corresponde a vários acessos
em execução, e por isso a mesma instrução pode ocupar as posições $A_1$ e $A_2$
de uma tripla. O mesmo vale para uma função chamada duas vezes pelo mesmo
fluxo. O Racebench anota defeitos exatamente com essa forma, de modo que uma
implementação que tratasse cada instrução como um acesso único perderia
defeitos sem emitir qualquer sinal de erro.

\subsection{Estágio 2, prioridade e mascaramento}

O segundo estágio descarta candidatos cujos fluxos comprovadamente não podem se
sobrepor. Um fluxo só é interrompido por outro de prioridade estritamente
maior, com a convenção de que número maior é prioridade maior. Fluxos de mesma
prioridade são tratados como capazes de se interromper mutuamente, pois nenhum
dos modelos formais da literatura cobre esse caso e essa é a leitura que não
perde defeitos.

O mascaramento é calculado por uma análise de fluxo de dados que determina, em
cada ponto do programa, quais interrupções estão desabilitadas. Uma interrupção
só conta como mascarada durante o intervalo $[A_1, A_2]$ se estiver mascarada
em todo ponto de todo caminho de $A_1$ a $A_2$ no fluxo local e se nenhum fluxo
capaz de executar dentro desse intervalo puder reabilitá-la. A segunda condição
é o que evita perder o defeito do caso \texttt{svp\_simple\_001\_001} descrito
na Seção~\ref{sec:interrupcoes}. Escritas em registradores que a configuração
não reconheça nunca contam como mascaramento.

Nenhum modelo de periodicidade é assumido, porque o Racebench declara que uma
interrupção pode ocorrer em qualquer ponto habilitado, quantas vezes for.

\subsection{Estágio 3, viabilidade com Z3}

O terceiro estágio aplica a análise de viabilidade de caminho de
\citeonline{intrace}. Como o estágio anterior, só descarta com prova, isto é,
com um veredicto de insatisfazibilidade. O resultado admite quatro valores,
contando o inconclusivo e o não executado, e um resultado inconclusivo, por
limite de tempo ou construção não suportada, nunca é tratado como decidido. O
candidato segue para a triagem carregando a informação de até onde o
solucionador chegou.

A posição desse estágio antes da triagem é fixa, por três motivos. Ele é um
filtro correto, que só descarta com prova. Reduz fortemente o volume, e os
números do IntRace dão a ordem de grandeza, com cerca de 208 candidatos por
programa caindo para cerca de 95 após o segundo estágio e para algo em torno de
13 após o terceiro \cite{intrace}, redução que é o que torna a triagem por \llm\
economicamente viável. E reproduz a arquitetura do LLift, em que o modelo
recebe apenas os casos que a execução simbólica não conseguiu decidir
\cite{llift}.

\subsection{O registro de contexto}

Para cada candidato sobrevivente a ferramenta emite um registro de contexto,
que é um documento estruturado com as evidências necessárias para julgá-lo.
\citeonline{llift} atribuem 5 de seus 13 falsos positivos a lacunas no
relatório entregue pelo analisador ao modelo, o que justifica tratar esse
registro como produto de primeira classe da análise, e não como subproduto. O
Quadro~\ref{quadro:registro} resume seus campos.

\begin{quadro}[htb]
  \caption{\label{quadro:registro}Campos do registro de contexto}
  \begin{tabular}{|l|L{10.4cm}|}
    \hline
    Campo & Conteúdo \\
    \hline
    padrão & qual dos quatro padrões do Quadro~\ref{quadro:padroes} \\
    \hline
    variável & nome, tipo, se é \texttt{volatile}, local da declaração \\
    \hline
    acessos & para cada um, leitura ou escrita, linha, função e fluxo \\
    \hline
    fluxos & ponto de entrada, prioridade e como o fluxo foi identificado \\
    \hline
    caminhos & sequência de chamadas do ponto de entrada até cada acesso \\
    \hline
    mascaramento & estado de cada interrupção em cada acesso e no intervalo \\
    \hline
    laços & laços que envolvem cada acesso, e se $A_1$ e $A_2$ são a mesma instrução \\
    \hline
    código & corpo completo das funções que contêm os acessos \\
    \hline
    solucionador & veredicto e, se inconclusivo, o motivo \\
    \hline
    proveniência & o que foi provado, o que veio da configuração, o que é desconhecido \\
    \hline
  \end{tabular}
  \fonte{Elaborado pelos autores.}
\end{quadro}

Duas ausências no registro são intencionais. Não há campo booleano de
mascaramento, e sim três valores, incluído o desconhecido, porque um booleano
obrigaria o desconhecido a ser representado como um dos outros dois. E não há fatiamento de programa
pré-calculado. Um fatiamento sobre uma variável global tende a abarcar quase
todo o programa, e um fatiamento incompleto leva o modelo a declarar com
confiança que um defeito real é impossível. O registro traz as camadas sempre
úteis, e o modelo pede o que falta, como descrito adiante.

\subsection{Triagem, duas perguntas e quatro grupos}

O modelo nunca escolhe diretamente onde o candidato fica na fila. Ele responde
a duas perguntas estreitas, a de se a intercalação pode acontecer e, caso
possa, a de se ela causa dano. Separar as perguntas evita que ``benigno'' se
converta silenciosamente em ``impossível'', que seria a confusão mais custosa
neste contexto. O grupo do candidato sai das duas respostas por uma
tabela fixa (Quadro~\ref{quadro:grupos}), em que qualquer dúvida sobe o
candidato na fila de leitura.

\begin{quadro}[htb]
  \caption{\label{quadro:grupos}Derivação do grupo de cada candidato, em ordem de leitura}
  \begin{tabular}{|c|L{3.4cm}|L{5.6cm}|L{2.8cm}|}
    \hline
    Ordem & Grupo & Viabilidade & Nocividade \\
    \hline
    1 & provavelmente real & possível & nociva \\
    \hline
    2 & incerto & incerta, ou possível com nocividade incerta & qualquer \\
    \hline
    3 & provavelmente benigno & possível & benigna \\
    \hline
    4 & provavelmente inviável & impossível, com o elemento que impede & não perguntada \\
    \hline
  \end{tabular}
  \fonte{Elaborado pelos autores.}
\end{quadro}

Todos os grupos aparecem no relatório. O grupo determina apenas a posição de
leitura, e é essa ordem que o \textit{Inspection Ratio} avalia.

\subsection{Desenho do \textit{prompt}}

O desenho segue os quatro componentes de \citeonline{llift}, adaptados ao
domínio. O primeiro é o ensino de regras do domínio por exemplos. O
conhecimento prévio do modelo vem sobretudo de \textit{threads} e o leva a
procurar travas e relações \textit{happens-before}, então o \textit{prompt}
ensina uma tabela de padrões de interrupção com seus veredictos, cobrindo a
seção crítica que abrange todo o intervalo, as duas seções separadas que deixam
uma lacuna no meio, a prioridade estritamente menor que impede a preempção, a
prioridade igual que não impede, e a mesma instrução num laço formando uma
tripla consigo mesma.

O segundo componente são os pedidos progressivos de contexto. O modelo pode
pedir, num formato fixo, a definição de uma função, o corpo de uma \isr, o
estado de mascaramento numa linha ou a expansão de uma macro. O registro desses
pedidos tem valor próprio, pois mede empiricamente que contexto o registro
deveria conter desde o início, questão que a literatura de fatiamento não
resolve para variáveis globais.

O terceiro é a decomposição da tarefa, com viabilidade e nocividade perguntadas
em conversas separadas. O quarto é a autovalidação enviesada para o recall,
com o modelo conferindo regras fixas antes de responder, entre elas a de que
mascaramento desconhecido conta como habilitado, a de que uma função cuja
definição não foi obtida pode acessar a variável, e a de que ``incerto'' é
sempre uma resposta segura.

A resposta final é estruturada, com a explicação antes do veredicto, como
recomendam \citeonline{iris}. O código analisado entra delimitado e declarado
como dado, conforme o RNF7.

\subsection{Reparo com testemunha}

Para cada candidato, o modelo pode propor uma correção no vocabulário de
interrupções enumerado por \citeonline{sdracer}, desabilitando e reabilitando a
interrupção em torno do trecho afetado, estendendo uma seção crítica existente
ou fundindo seções adjacentes. A correção precisa tornar atômico todo o
intervalo $[A_1, A_2]$, e não apenas proteger um dos acessos isoladamente.

Inspirada na validação por \textit{exploits} de \citeonline{sastgenius}, toda
proposta exige uma testemunha, que diz qual \isr\ dispara, em que ponto, em que
ordem ocorrem os acessos e que estado resulta diferente de qualquer execução em
série. A testemunha é conferida mecanicamente contra o registro de contexto antes de a
correção ser aceita, o que impede a aceitação de uma correção plausível, porém
injustificada.

\subsection{Reverificação}

Uma correção é verificada em duas etapas que falham em direções opostas.
Primeiro a análise estática é executada de novo sobre o programa corrigido.
Como seus filtros só descartam com prova, o desaparecimento do candidato
certifica, sob as premissas declaradas, que o defeito foi eliminado, e a nova análise
completa é a única forma de detectar regressões, como uma seção crítica
estendida que cria outro defeito. Se o candidato persistir, o resultado é
inconclusivo, porque uma análise superaproximada pode continuar reportando um
defeito já corrigido. Os casos inconclusivos vão então para um verificador
limitado, que codifica a violação como uma asserção. Por fim, como observam
\citeonline{sdracer}, verifica-se se a correção alterou o fluxo de controle ou
a latência, que em código de interrupção constitui propriedade de correção, e
não mero detalhe de desempenho.

\section{Testes e Avaliação}
\label{sec:avaliacao}

\subsection{Perguntas de pesquisa}

A avaliação responde a quatro perguntas. A primeira é se a etapa estática
encontra todos os defeitos anotados no Racebench, ou seja, se atinge recall de
100\%. A segunda é se a triagem reduz o esforço de revisão sem que nenhum
defeito real seja classificado como inviável. A terceira é qual a contribuição
de cada técnica de \textit{prompt}, respondida por uma ablação no formato de
\citeonline{llift}. A quarta é se o modelo propõe correções corretas para os
defeitos encontrados.

\subsection{Unidade de contagem e leitura do gabarito}

Toda taxa de recall divide por um número de defeitos, e a
Seção~\ref{sec:benchmarks-conceito} mostrou que a literatura usa números
diferentes para os mesmos programas sem justificar a escolha. Este trabalho
adota uma tripla anotada como um defeito. A Tabela~\ref{tab:unidades} mostra
quanto a escolha muda o denominador.

\begin{table}[htb]
  \centering
  \caption{Contagem do gabarito do Racebench conforme a unidade}
  \label{tab:unidades}
  \begin{tabular}{|l|c|c|}
    \hline
    Unidade & Defeitos & Armadilhas \\
    \hline
    por instância de tripla & 48 & 38 \\
    \hline
    por variável em cada caso & 33 & 29 \\
    \hline
    por caso & 31 & 31 \\
    \hline
  \end{tabular}
  \fonte{Elaborada pelos autores a partir das anotações de \citeonline{racebench}.}
\end{table}

A unidade por instância é a mais fina e, por isso, a mais difícil de inflar. É
também a usada pelo BMC4AV, o que preserva a comparação, e as unidades mais
grossas escondem estrutura, já que um mesmo caso chega a ter quatro defeitos
distintos numa só variável.

As anotações usam ao menos quatro gramáticas incompatíveis e contêm erros, de
modo que um leitor estrito encontra apenas 28 defeitos sem acusar nenhuma
falha. O leitor de gabarito precisa portanto tolerar as variações, e sua
contagem é validada contra uma contagem manual dos casos mais difíceis antes de
qualquer medição. Doze anotações são corrigidas por uma errata em que cada
correção traz sua evidência, sendo seis números de linha, cinco deles num único
caso cujas anotações apontam para linhas sem nenhum acesso, e seis tipos de
acesso trocados. Nenhuma correção altera as contagens. Sem elas, uma comparação
exata com o texto distribuído daria recall zero para um defeito que a
ferramenta encontra na linha certa.

\subsection{Regra de casamento}

Um candidato reportado corresponde a uma tripla anotada quando está no mesmo
programa, trata da mesma posição de memória e tem as mesmas linhas ordenadas
$(A_1, B, A_2)$ que a anotação corrigida. Três escolhas são deliberadas. O tipo
de acesso não precisa coincidir, porque as linhas já distinguem todas as 86
anotações sem colisão e 20 acessos anotados estão em linhas que leem e escrevem
a variável. O nome da variável é comparado depois de resolver apelidos, pois 12
anotações nomeiam a posição por um ponteiro, um elemento de vetor ou um campo
de estrutura. E a comparação é feita sobre as anotações corrigidas, cuja errata
é publicada junto com os resultados.

\subsection{Conjunto de triagem}
\label{sec:conjunto-triagem}

A etapa com \llm\ é desenvolvida e medida, até a integração, sobre 20 registros
de contexto montados à mão a partir do Racebench, sendo 10 defeitos anotados e
10 armadilhas. Eles estão organizados de modo a formar seis duplas
adversariais, isto é, um defeito e uma armadilha que diferem num único detalhe,
como o estado de mascaramento. As duplas existem para que o modelo não possa
acertar com base apenas na aparência do código. Um caso cujas anotações descrevem uma versão do programa
diferente da distribuída foi excluído do conjunto.

A suíte de 18 programas reais entra na etapa final. Seu gabarito registra, por
variável, 45 violações verdadeiras e 14 falsas, o que também a torna um
conjunto de triagem.

\subsection{Métricas}

O portão de recall exige que todos os defeitos anotados cheguem ao relatório e
que nenhum caia no grupo de menor prioridade de leitura. Uma falha aqui é
defeito de projeto, não parâmetro a ajustar. A rejeição de armadilhas mede a
fração das armadilhas classificadas nos grupos baixos, e é a medida de precisão
da triagem que, ao contrário da precisão, não melhora descartando candidatos. O
\textit{Inspection Ratio} mede a fração da lista ordenada que um revisor
precisa ler até encontrar todos os defeitos \cite{reducing-fp}, com empates
resolvidos pelo pior caso, colocando o defeito real por último. A consistência
mede a concordância entre execuções repetidas do mesmo candidato, e
\citeonline{llift} mostram que decomposição e autovalidação a melhoram
independentemente da acurácia. As taxas de reparo são duas, a sintática, em que
um analisador sintático aceita a correção, e a lógica, em que a correção
equivale a uma escrita à mão, como em \citeonline{skipanalyzer}, seguidas da
reverificação.

No Racebench apenas a viabilidade entra nessas métricas. O gabarito anota fatos
do programa, e muitos de seus defeitos leem a variável compartilhada para uma
variável local que nunca é usada. Um modelo que classifica esses casos como
benignos está raciocinando corretamente sobre o código, e pontuar essa resposta
contra o gabarito puniria o acerto. A pergunta de nocividade continua sendo
feita e registrada, e é avaliada na suíte de programas reais, na qual o código tem
efeito observável.

\subsection{Protocolo experimental}

A ablação acrescenta um componente de \textit{prompt} por linha, sempre sobre
os mesmos candidatos, partindo do \textit{prompt} simples e somando as regras
do domínio, os pedidos progressivos de contexto, a decomposição da tarefa e a
autovalidação. Como os pedidos de contexto só trazem informação nova quando
respondidos a partir do código real, a linha correspondente é medida depois da
integração com o resolvedor da etapa estática.

Na ablação do LLift o \textit{prompt} simples já perdia a maioria dos defeitos,
e o recall era a variável que se movia. Se aqui o recall partir de 100\%, as
variáveis dependentes passam a ser a rejeição de armadilhas e o
\textit{Inspection Ratio}, e é assim que a tabela final deve ser lida.

Todo número é registrado com o modelo, a data e a configuração de
\textit{prompt} que o produziram, conforme o RNF6, e uma linha da ablação só é
considerada medida se cobrir todos os candidatos previstos.

\ifparcial\else
\subsection{Resultados}
\label{sec:resultados}

\pendente{Recall da etapa estática, tabela de ablação da triagem, rejeição de
armadilhas, \textit{Inspection Ratio}, taxas de reparo, custo por candidato e
resultados na suíte de programas reais.}
\fi

\subsection{Ameaças à validade}
\label{sec:ameacas}

Vinte candidatos e uma amostra por candidato tornam diferenças de um ou dois
candidatos indistinguíveis de ruído, e é a medição de consistência com várias
amostras que permite afirmar uma diferença. Quando o \textit{prompt} é
corrigido depois de observar erros nos mesmos 20 casos, o número seguinte é
otimista, e por isso as versões anteriores e posteriores são reportadas lado a
lado. A errata corrige o que pode ser verificado no código, mas o rótulo de um
caso cujas anotações descrevem outra versão do programa não pode ser corrigido
sem decidir de novo o gabarito.

Duas ameaças são estruturais. Depois da integração a população muda, pois o
modelo passa a receber apenas o que o solucionador não decidiu, que são os
casos mais difíceis, de modo que os números obtidos sobre os registros montados
à mão não se transferem automaticamente. E todos os programas avaliados são de
arquivo único, o que deixa a entrada com vários arquivos como capacidade
exercida fora do gabarito. Por fim, a exclusão de reentrância da
Seção~\ref{sec:premissas} é uma perda potencial de recall registrada, não
medida.
